\documentclass[11pt,a4paper]{letter}
 
\usepackage[utf8]{inputenc} % for unicode typing
\usepackage[ngerman]{babel} % Neue Deutsch-Rechtschreibung
\usepackage[T1]{fontenc} % fonts encoding
\usepackage[]{newtxtext} % arial font

\usepackage{graphicx} % insert image and picture
\usepackage[left=2.5cm,
            right=2cm,
            top=2.5cm,
            bottom=2cm]{geometry} % define pagelayout

% using microtype package, so to text could be better format
\usepackage[activate={true,nocompatibility},
            final,
            tracking=true,
            kerning=true,
            spacing=true,
            factor=1100,
            stretch=10,
            shrink=10]{microtype}

\usepackage{bookmark} % create a custom bookmark

\setlength{\parskip}{1em} % set the space between two paragraph
\setlength{\parindent}{0pt} % set the indent space of paragraph
\setcounter{secnumdepth}{2} % set the depth of the document, section

\usepackage{xargs} % use argument in new command
\usepackage[table,dvipsnames]{xcolor} % add color support for text

% TODO notes
\usepackage[colorinlistoftodos,prependcaption,textsize=footnotesize]{todonotes} % write TODO
\newcommandx{\unsure}[2][1=]{\todo[linecolor=red, backgroundcolor=red!25, bordercolor=red, #1]{#2}}
\newcommandx{\change}[2][1=]{\todo[linecolor=blue, backgroundcolor=blue!25, bordercolor=blue, #1]{#2}}
\newcommandx{\info}[2][1=]{\todo[linecolor=OliveGreen, backgroundcolor=OliveGreen!25, bordercolor=OliveGreen, #1]{#2}}
\newcommandx{\improvement}[2][1=]{\todo[linecolor=Plum, backgroundcolor=Plum!25, bordercolor=Plum, #1]{#2}}

% macros, define in documents
\newcommand{\FirmaName}{CTE Entwicklungen GmbH}
\newcommand{\FirmaAdresseLineOne}{Max-von-Laue-Straße 5}
\newcommand{\FirmaAdresseLineTwo}{86156 Augsburg}

\newcommand{\JobName}{Konstruktion und Entwicklung}
\newcommand{\JobBereich}{}
\newcommand{\JobRefNummer}{WO 71030}

\newcommand{\AnsprechpartnerVoll}{Frau Ellena Domes}
\newcommand{\Ansprechpartner}{Maier}
\newcommand{\Ansprechpartnerin}{Domes}

\newcommand{\Beginntermin}{Oktober. 2017}

\begin{document}
% Manual import bookmark
\bookmark[page=1, level=-2]{Anschreiben}
% pagestyle empty means no page number
\pagestyle{empty}

% Sender
\begin{flushleft}
    \textbf{Ngoc Minh Dao}\\
    Obermichelbacher Str. 2\\
    90587 Obermichelbach\\
    E-Mail: ngocminhdao88@gmail.com\\
    Tel.: (+49) 176 74561589
\end{flushleft}

% Empfänger
\begin{flushleft}
    \textbf{\FirmaName}\\
    \AnsprechpartnerVoll\\
    \FirmaAdresseLineOne\\
    \FirmaAdresseLineTwo\\
\end{flushleft}

% Today
\begin{flushright}
    Obermichelbach, \today
\end{flushright}

%\textbf{Bewerbung um eine Abschlussarbeit zum Thema: \JobName~(\JobRefNummer) ab \Beginntermin}
\textbf{Bewerbung um eine Abschlussarbeit ab \Beginntermin~im Bereich \JobName}\\
%Ref: \JobRefNummer
\vspace{1em}

%Sehr geehrte Damen und Herren,
%Sehr geehrter Herr \Ansprechpartner,
Sehr geehrte Frau \Ansprechpartnerin,

%als europaweit operierender Dienstleister der Fahrzeug- und Fahrzeugzulieferindustrie suchen Sie für den weiteren Ausbau Ihres Unternehmens einen hoch motivierten Student, ich suche eine herausordernde Aufgabe für den Abschluss meines Maschinenbau Studiums.
als Student des Studiengangs Maschinenbau mit der Vertiefungsrichtung Fahrzeugtechnik ist, sind Sie mir als der führende Hersteller für Landmaschinen von großem Interesse.

Schon während meines Studiums des Maschinenbaus konnte ich praktische technische Erfahrung gesammelt.
Als HiWi unterstützte ich das Team des Instituts für Maschinenkonstruktion und Tribologie bei der Konstruktionsaufgabe, dem Aufbau sowie der Durchführung von Versuchen für verschiedene Forschungsprojekte.

%\improvement[inline]{Passen die Aufgaben während des Praktikums mit dem Thema?}

Der Schwerpunkt legte ich in meinem Studium auf die Entwicklung von Getrieben und Antriebssträngen sowie die Anwendung von Sensoren in derartigen Systemen.
Im Rahmen meiner Projektarbeit habe ich ein Prüfkonzept zur Untersuchung der Dichtheit von rotierenden Radialwellendichtringen entwickelt.
Neben dem Studium erweiterte ich meine technischen Kenntnisse dadurch, dass ich an Online-Lehrangeboten, Seminaren und Tutorien teilnahm.
Außeruniversitär war ich ein Mitglied in einer akademischen Gruppe für Kraftfahrwesen, wo ich in meiner Freizeit anderen Kommilitonen bei der Wartung bzw. der Reparatur von Fahrzeugen unterstützte.

%\improvement[inline]{Welchen konkreten Schwerpunkt möchte ich mit der Abschlussarbeit bearbeiten?}

Weitere Erfahrung sammelte ich in meinem Praktikum bei Schaeffler Gruppe.
Hier unterstützte ich das Getriebe-Anwendungstechnik Team mit der Berechnung, der Schadensanalyse von Lagern, sowie dem Empfang und dem Büromanagement.
Hierbei konnte ich meine Teamfähigkeit unter Beweis stellen.

Aus der Kununu-Webseite wurde ich auch erfahren, dass Sie ein fairer Arbeitgeber sind.
Für Ihre Mitarbeiter bieten Sie nicht nur ein freundliches Arbeitsklima, sondern auch anspruchsvolle, interessante Aufgaben und gute Weiterbildungsmöglichkeiten an.
Dies verstärkt meine Wahl, bei Ihnen meine Abschlussarbeit zu schreiben.

Ich spreche Deutsch, Englisch und habe gute Kenntnisse der Tools Creo, Ansys, Matlab und MS Office.

Für diese Stelle kann ich sofort antreten.
Gerne überzeuge ich Sie in einem persönlichen Gespräch.

Mit freundlichen Grüßen

%\includegraphics[scale=0.2]{./signaturen/NgocMinhDao-Signatur.pdf}\\
Ngoc Minh Dao

Anlagen
\end{document}
